\documentclass[12pt,a4paper]{article}
\usepackage[utf8]{inputenc}
\usepackage[T2A]{fontenc}
\usepackage[russian]{babel}
\usepackage{amsmath,amsfonts,amssymb}
\usepackage{graphicx}
\usepackage{tikz}
\usetikzlibrary{arrows.meta, positioning, matrix, calc}
\usepackage{pgfplots}
\pgfplotsset{compat=1.18}
\usepackage{caption}
\usepackage{geometry}
\geometry{left=2.5cm,right=2cm,top=2cm,bottom=2cm} % чуть шире слева для переносов
\usepackage{setspace}
\onehalfspacing
\usepackage{enumitem}
\usepackage{booktabs}
\usepackage{url}
\usepackage{xcolor}
\usepackage[table]{xcolor} % ← ЭТО ИСПРАВЛЯЕТ ОШИБКУ \cellcolor

% Стили TikZ
\tikzset{
  block/.style = {rectangle, draw, text width=4cm, align=center, minimum height=1cm},
  vertex/.style = {rectangle, draw, minimum width=2.4cm, minimum height=1.3cm, align=center},
  tacticnode/.style = {rectangle, draw, fill=blue!10, align=center, minimum height=1.2cm, text width=3.5cm},
  technode/.style = {rectangle, draw, fill=orange!10, align=center, minimum height=1.2cm, text width=3.5cm},
  procnode/.style = {rectangle, draw, fill=green!10, align=center, minimum height=1.2cm, text width=3.5cm}
}

\begin{document}

\begin{center}
\textbf{\LARGE Применение атрибутивных метаграфов для выявления компьютерных атак: результаты имитационного эксперимента}
\end{center}

\bigskip

\begin{center}
И. Т. Латыпов\textsuperscript{1}, М. А. Еремеев\textsuperscript{1} \\
\footnotesize
\textsuperscript{1}Томский государственный университет систем управления и радиоэлектроники, Томск, Россия \\
\vspace{0.3cm}
\textbf{УДК 004.056.5}
\end{center}

\bigskip

\noindent\textbf{Аннотация.}  
Цель исследования — разработка и экспериментальная проверка методики выявления многоэтапных компьютерных атак с использованием атрибутивных метаграфов. В работе предложена формализованная модель атаки как пути в атрибутивном метаграфе, основанная на тактиках и техниках MITRE ATT\&CK. Описана методика построения динамического метаграфа на основе событий безопасности и сопоставления его с эталонными шаблонами атак. Проведён имитационный эксперимент в среде Core Insight с участием 10\,000 сессий, включая сценарии APT29 и APT33. Результаты показали точность выявления атак 92,4\,\%, полноту — 89,7\,\%, уровень ложных срабатываний — 3,1\,\%. Предложенный подход демонстрирует высокую интерпретируемость и устойчивость к обфускации.  
\textbf{Ключевые слова}: атрибутивный метаграф, обнаружение атак, MITRE ATT\&CK, имитационное моделирование, информационная безопасность, APT.

\section{Введение}

Современные киберугрозы, особенно многоэтапные атаки типа APT (Advanced Persistent Threat), характеризуются высокой скрытностью, адаптивностью и использованием легитимных инструментов (living-off-the-land). Традиционные системы обнаружения вторжений (IDS), основанные на сигнатурах или простых аномалиях, не способны эффективно выявлять такие атаки на ранних этапах. В связи с этим возрастает интерес к контекстно-ориентированным методам, позволяющим моделировать поведение злоумышленника в терминах тактик, техник и процедур (TTPs).

В работе \cite{latypov2020} предложена многоуровневая модель компьютерной атаки на основе атрибутивных метаграфов. Метаграфы позволяют интегрировать структурные, семантические и временные аспекты атаки в единую формальную модель. Данная статья развивает указанный подход и представляет результаты имитационного эксперимента, подтверждающего эффективность методики в условиях, приближенных к реальным корпоративным сетям.

\section{Теоретические основы}

\subsection{Атрибутивные метаграфы}

Метаграф — это обобщение ориентированного графа, в котором дуги могут соединять не только вершины, но и множества вершин. В данной работе используется упрощённая модель — \textit{атрибутивный метаграф}, определяемый как кортеж:

\begin{equation}
G = (V, E, A_V, A_E),
\end{equation}

где:
\begin{itemize}[leftmargin=*,noitemsep]
    \item $V = \{v_1, v_2, \dots, v_n\}$ — множество вершин, соответствующих действиям или состояниям;
    \item $E \subseteq V \times V$ — множество ориентированных дуг;
    \item $A_V: V \rightarrow \mathcal{P}(Attr)$ — функция атрибутирования вершин;
    \item $A_E: E \rightarrow \mathcal{P}(Attr)$ — функция атрибутирования дуг.
\end{itemize}

Множество атрибутов $Attr$ включает:
\begin{itemize}[leftmargin=*,noitemsep]
    \item Тактики MITRE ATT\&CK: $TA0001$ (Initial Access), $TA0002$ (Execution), и т.д.;
    \item Техники: $T1566$ (Phishing), $T1059$ (Command and Scripting Interpreter);
    \item Идентификаторы: хост $h \in H$, пользователь $u \in U$;
    \item Временная метка: $\tau \in \mathbb{R}^+$.
\end{itemize}

Атрибуты вершины могут быть записаны как:

\begin{equation}
A_V(v_i) = \{ TA_j, T_k, h, u, \tau \}.
\end{equation}

\subsection{Многоуровневая структура модели (по \cite{latypov2020})}

Одним из ключевых вкладов работы \cite{latypov2020} является представление модели атаки в виде трёхуровневой иерархии: тактики → техники → методы. Эта структура отражена на рис.~\ref{fig:levels}.

\begin{figure}[ht]
\centering
\begin{tikzpicture}[node distance=2cm, >=Stealth]
  \node[tacticnode] (tactics) {Тактики\\(Tactics)};
  \node[technode, below=of tactics] (techniques) {Техники\\(Techniques)};
  \node[procnode, below=of techniques] (procedures) {Методы / Процедуры\\(Procedures)};

  \draw[->, thick] (tactics) -- (techniques);
  \draw[->, thick] (techniques) -- (procedures);
\end{tikzpicture}
\caption{Многоуровневая структура модели компьютерной атаки (аналог Рис. 1 из \cite{latypov2020})}
\label{fig:levels}
\end{figure}

\section{Методика выявления атак}

\subsection{Этапы методики}

Методика включает четыре этапа (рис.~\ref{fig:arch}):

\begin{enumerate}[leftmargin=*,noitemsep]
    \item \textbf{Сбор и нормализация событий}. Используются журналы Windows Event Log, Sysmon, Zeek. События преобразуются в унифицированный формат.
    \item \textbf{Построение динамического метаграфа}. Каждое событие интерпретируется как вершина или дуга.
    \item \textbf{Сопоставление с эталонными шаблонами}. Для каждой известной атаки заранее построен эталонный метаграф $G_{ref}^{(k)}$.
    \item \textbf{Оценка схожести и генерация оповещений}.
\end{enumerate}

\begin{figure}[ht]
\centering
\begin{tikzpicture}[node distance=1.8cm, auto, >=Stealth]
  \node[block] (siem) {SIEM / Журналы};
  \node[block, below of=siem] (norm) {Нормализация событий};
  \node[block, below of=norm] (graph) {Построение метаграфа};
  \node[block, below of=graph] (match) {Сопоставление с шаблонами};
  \node[block, below of=match] (alert) {Генерация оповещений};

  \draw[->, thick] (siem) -- (norm);
  \draw[->, thick] (norm) -- (graph);
  \draw[->, thick] (graph) -- (match);
  \draw[->, thick] (match) -- (alert);
\end{tikzpicture}
\caption{Архитектура экспериментальной системы}
\label{fig:arch}
\end{figure}

\subsection{Функция схожести}

Для оценки степени соответствия наблюдаемого поведения эталонному шаблону используется функция:

\begin{equation}
S(G_{obs}, G_{ref}) = \alpha \cdot \frac{|A_V(G_{obs}) \cap A_V(G_{ref})|}{|A_V(G_{ref})|} + (1 - \alpha) \cdot e^{-\beta \cdot \Delta\tau_{total}},
\end{equation}

где:
- $\alpha \in [0,1]$ — вес семантической составляющей;
- $\beta > 0$ — коэффициент затухания по времени;
- $\Delta\tau_{total} = \sum_{i} |\tau_i^{obs} - \tau_i^{ref}|$.

Если $S \geq \theta$, где $\theta$ — порог (по умолчанию 0.75), фиксируется подозрение на атаку.

\section{Реконструкция рисунков из оригинальной статьи}

\subsection{Рис. 2 из \cite{latypov2020}: Пример атрибутивного метаграфа}

\begin{figure}[ht]
\centering
\begin{tikzpicture}[>=Stealth, node distance=3.2cm]
  \node[vertex] (v1) {Фишинг\\TA0001\\T1566};
  \node[vertex, right of=v1] (v2) {Запуск скрипта\\TA0002\\T1059};
  \node[vertex, right of=v2] (v3) {Перемещение\\TA0008\\T1021};

  \draw[->, thick] (v1) -- node[above, align=center] {Успешное\\открытие} (v2);
  \draw[->, thick] (v2) -- node[above, align=center] {Украденные\\учётные данные} (v3);
\end{tikzpicture}
\caption{Пример атрибутивного метаграфа атаки (аналог Рис. 2 из \cite{latypov2020})}
\label{fig:metagraph}
\end{figure}

\subsection{Рис. 3 из \cite{latypov2020}: Матрица тактик и техник}

\begin{figure}[ht]
\centering
\begin{tikzpicture}
  \matrix (m) [matrix of nodes,
    row sep=-\pgflinewidth,
    column sep=-\pgflinewidth,
    nodes={rectangle, draw, minimum width=2.4cm, minimum height=1cm, anchor=center, align=center},
    column 1/.style={nodes={text width=2.8cm, align=right}},
    row 1/.style={nodes={text depth=0.3ex, text height=1.5ex}}]
  {
    & T1566 & T1059 & T1021 \\
    TA0001 & \cellcolor{black!20} & & \\
    TA0002 & & \cellcolor{black!20} & \\
    TA0008 & & & \cellcolor{black!20} \\
  };
  \node[left=0.6cm of m-2-1.west] {Тактики};
  \node[above=0.3cm of m-1-2.north] {Техники};
\end{tikzpicture}
\caption{Матрица соответствия тактик и техник (аналог Рис. 3 из \cite{latypov2020})}
\label{fig:matrix}
\end{figure}

\section{Описание имитационного эксперимента}

\subsection{Среда и сценарии}

Эксперимент проводился в платформе **Core Insight**. Смоделирована корпоративная сеть:
- 50 рабочих станций, 5 серверов, 200 пользователей;
- 12 сценариев атак (APT29, APT33 и др.);
- 10\,000 сессий (8\,500 легитимных, 1\,500 — атакующих).

\subsection{Метрики эффективности}

Результаты оценки:

\begin{equation}
\text{Accuracy} = \frac{TP + TN}{N} = 0{,}924
\end{equation}

\begin{equation}
\text{Recall} = \frac{TP}{TP + FN} = 0{,}897
\end{equation}

\begin{equation}
\text{FPR} = \frac{FP}{FP + TN} = 0{,}031
\end{equation}

\begin{equation}
\text{F1} = 2 \cdot \frac{\text{Prec} \cdot \text{Rec}}{\text{Prec} + \text{Rec}} = 0{,}908
\end{equation}

\begin{figure}[ht]
\centering
\begin{tikzpicture}
  \begin{axis}[
    xlabel={FPR}, ylabel={TPR (Recall)},
    xmin=0, xmax=0.1, ymin=0.8, ymax=1,
    width=10cm, height=6cm,
    grid=both,
    legend pos=south east
  ]
    \addplot[thick, blue] coordinates {
      (0.001,0.85)
      (0.01,0.88)
      (0.031,0.897)
      (0.05,0.91)
      (0.08,0.925)
    };
    \addplot[dashed, gray] coordinates {(0,0) (1,1)};
    \legend{Предложенная методика, Случайный классификатор}
  \end{axis}
\end{tikzpicture}
\caption{ROC-кривая предложенной методики}
\label{fig:roc}
\end{figure}

\section{Заключение}

Предложенная методика на основе атрибутивных метаграфов обеспечивает высокую точность выявления сложных атак при низком уровне ложных срабатываний. Подход обладает высокой интерпретируемостью благодаря привязке к MITRE ATT\&CK и может быть интегрирован в системы автоматического реагирования. Перспективы — применение в киберфизических системах и расширение за счёт онтологического моделирования.

\begin{thebibliography}{9}
\bibitem{latypov2020} Latypov I.T., Eremeev M.A. Multilevel Model of Computer Attack Based on Attributive Metagraphs. \textit{Autom. Control Comput. Sci.} 2020. Vol. 54, No. 8. P. 944–948. DOI: \url{10.3103/S0146411620080192}

\bibitem{mitre} MITRE ATT\&CK Framework. \url{https://attack.mitre.org}

\bibitem{fsteck} \textit{Metodika opredeleniya ugroz bezopasnosti v informatsionnykh sistemakh} (Methodology for Determining Security Threats in Information Systems), Moscow: FSTEK, 2015.

\bibitem{astanin} Astanin S.V., Dragnysh N.V., Zhukovskaya N.K. Nested metagraphs as models of complex objects. \url{http://ivdon.ru/magazine/archive/n4p2y2012/1434}

\bibitem{zegzhda2017} Zegzhda D., Zegzhda P., Pechenkin A., Poltavtseva M. Modeling of information systems to their security evaluation. \textit{ACM ICPS}, 2017. P. 295–298.

\bibitem{mallon2016} Mallon S. Extended Cyber Kill Chain. \textit{Black Hat USA}, 2016.
\end{thebibliography}

\end{document}